La temperatura en una placa metálica está dada por $T(x, y) = 100 - 3x^2 - 5y^2$ (en $^\circ$C), con $x$ e $y$ en metros.
\begin{enumerate}
    \item Encuentre el gradiente de $T$ en el punto $(1, 2)$.
    \item En ese punto, ¿en qué dirección debe moverse un punto de la placa para que la temperatura suba con mayor rapidez? ¿Cuál es esa tasa máxima de cambio?
    \item Calcule la derivada direccional en la dirección que forma $\pi/4$ con el eje $x$.
\end{enumerate}
